\section{Imitación jerárquica y descubrimiento de habilidades}

\subsection{Extracción de habilidades a partir de demostraciones}

El aprendizaje por imitación ya no depende de la reward, sino que infiere la skill segmentation a partir de la demo. Entre los enfoques habituales están sliding window clustering, change point detection y el skill-specific autoencoder.

\begin{importantbox}{Pipeline de descubrimiento de habilidades}
\begin{enumerate}
    \item Etiquetado: dividir la demo en candidate skills;
    \item Entrenar la skill policy: usar el segment como policy context;
    \item Entrenar el scheduler: el nivel alto aprende a predecir la skill sequence;
    \item Deployment: vincular la skill library a un planificador LLN/LLM.
\end{enumerate}
\end{importantbox}

\subsection{Gobernanza de los datos de imitación}

\begin{knowledgebox}{Checklist para gobernar los datos de imitación}
\begin{itemize}
    \item Demo provenance: registrar el id del video o la trayectoria de origen;
    \item Skill annotation: anotar los límites de las subtareas y exportarlos a `labels.json`;
    \item Diversity filter: eliminar los fragmentos de demostración excesivamente repetidos;
    \item Success flag: marcar en cada demostración si se completó, para que la use la low-layer reward.
\end{itemize}
\end{knowledgebox}

\subsection{VQ-VAE y espacio latente discreto}

El VQ-VAE se encarga de comprimir trayectorias variable-length en discrete codebook entries; cada entry corresponde a una habilidad. La política de nivel alto solo necesita predecir el codebook index, y el decoder de nivel bajo puede «decodificar» la secuencia de acciones necesaria.

\begin{figure}[H]
\centering
\includegraphics[width=0.85\textwidth]{slides-images/slide-25.jpg}
\caption{Esquema de la arquitectura de Skill Discovery (slide-25).}
\end{figure}

\begin{warningbox}{Collapse de la variable latente}
Si el codebook size es demasiado pequeño, el VQ-VAE reutiliza una y otra vez la misma habilidad, lo que deja a la política de nivel alto sin diversidad. Las soluciones habituales son añadir una commitment loss o usar EM-style clustering.
\end{warningbox}

\subsection{El lenguaje como habilidad y como comando}

Un LLM puede servir como interface de lenguaje para la policy de nivel alto; el ponente mencionó que SayCan/Inner Monologue, mediante el pipeline language → intent → skills, aumentan la interpretabilidad del ejecutor de nivel bajo.

\begin{knowledgebox}{Puntos clave para gobernar las instrucciones en lenguaje}
\begin{itemize}
    \item El prompt log registra la language input + la expected skill;
    \item Usar un parser para convertir el language en un discrete skill id;
    \item La retroalimentación del nivel bajo (por ejemplo, el success flag) debe escribirse de vuelta en el prompt log, para formar un governance loop.
\end{itemize}
\end{knowledgebox}

\subsection{Resumen de la sección}

El módulo de imitación permite descubrir automáticamente una biblioteca de habilidades a partir de demostraciones; junto con una interfaz de lenguaje/LMM, logra una planificación jerárquica que las personas pueden interpretar.
